%% Chapter 7 - Discussion
%% Sections: 7.1 Interpretation of findings; 7.2 Comparison with literature
%% baselines and existing solutions; 7.3 Research contribution; 7.4 Product
%% contribution; 7.5 Limitations and threats to validity; 7.6 Implications.

\chapter{Discussion}\label{ch:discussion}

This chapter interprets the measurements reported in
Chapter~\ref{ch:results} against the research questions of
Section~\ref{sec:research-questions}, compares the guardrail stack with the
surveyed baselines of Chapter~\ref{ch:literature-review}, and separates the
research from the product contribution. It then states the verified limitations
and threats to validity and draws the practical and research implications.

\section{Interpretation of Findings}\label{sec:discussion-interpretation}

The guarded pipeline answers the central question affirmatively: a small,
offline, unattended language model can be made reliable enough to gate an
automated security copilot, provided its output never reaches downstream
tooling before it satisfies a conjunction of deterministic predicates. The
measured basis for this claim is the experiment record reported in
Chapter~\ref{ch:results}, whose headline deltas are considered under each
research question in turn.

\subsection*{RQ1: pipeline value}

The raw model is not merely imprecise---it is unusable as an intent
classifier. Experiment~1
(Section~\ref{sec:exp1}) shows a catch rate of $0\%$ on both backends:
a single-shot schema probe never happens to produce a flaw-free intent for
the nine-category taxonomy, and the false-positive rate compounds the
problem---32\% for Gemma and 31\% for Qwen of well-formed records are rejected
or misclassified. The guarded pipeline converts this into $D_{\mathrm{Acc}}$
86\%/82\%, catch 72\%/64\% and FP 14\%/15\% (Gemma/Qwen,
Table~\ref{tab:exp1}). The accuracy gain is therefore not the headline
result: it is the catch rate, which the guards raise from zero to between two
thirds and three quarters of all flawed records while simultaneously halving
the false-positive rate. The most useful single reading of RQ1 is the
asymmetry between the two movements---the deterministic stack adds the 
ability to detect flaws, and reduces spurious rejection at the same time,
rather than trading one for the other.

\subsection*{RQ2: layer contribution}

Experiment~3 (Table~\ref{tab:ablation}) identifies the validation stack as the
decisive component: removing it entirely collapses catch from 72\% to 42\% for
Gemma and 64\% to 40\% for Qwen ($\Delta$Catch $-30$ and $-24$). Individually,
the hallucination validators are the largest single-layer contributors to catch
($-9$ Gemma, $-14$ Qwen), which is consistent with the dominant flaw classes in
the synthetic population: fabricated parameters and CVEs are the most frequent
flaw kinds. The design insight from Experiment~4
(Section~\ref{sec:exp4}) is that per-class catch accords with determinism:
the classes that expose a syntactic or enumerable signal (RBAC violations,
adversarial phrasing, fabricated parameters and CVE identifiers) are caught at
78--100\%, whereas the classes that require meaning resolution
(contradictory, out-of-scope, ambiguous) remain at 40--59\%. RQ2 is therefore
answered both at the layer level and at the class level: the guard layers that
contribute most are exactly the ones that encode decidable predicates, and the
residual risk concentrates in the predicates that are not decidable.

The negative $\Delta$Catch entries in the left column of
Table~\ref{tab:ablation}---removing the ablation-row guard apparently raising
catch---are read, as Section~\ref{sec:data-analysis} commits, not as a
deficiency of running the counterfactual treatment, but as a removal-adversity
finding: the default configuration trades a small portion of catch against a
large preservation of usability (false positives and accuracy). This is
expected under the framework of Chapter~\ref{ch:literature-review} and does not
contradict the safety invariant; it is artefact behaviour of the committed
single-run protocol rather than evidence about the production setting.

\subsection*{RQ3: generalisability}

The catch-rate ordering of the ablation rows is broadly consistent across the
two backends (Experiment~3), so RQ3 is answered affirmatively for the safety
dimension: deterministic guardrailing, not model choice, is the primary driver
of catch-rate performance, provided the trust boundary remains closed. The
condition on that conclusion is the accuracy dimension. The accuracy
contribution of individual layers is backend-specific---the scope guard
dominates accuracy for Gemma, the hallucination validator for Qwen---so the 
generality claim must be restricted to the ordering of catch-rate effects
(Section~\ref{sec:data-analysis}) and explicitly must not be extended to a
claim that any single layer is required in the same rank order to preserve
accuracy across backends. The measured safety therefore generalises; the
measured usability-tuning does not.

\section{Comparison with Literature Baselines and Existing Solutions}\label{sec:discussion-comparison}

Against the existing-systems survey
(Table~\ref{tab:existing-systems}), the IRE occupies the combination that the
review found empty: a dedicated intent-recognition layer that is
machine-consumed, schema-checked, scope- and authorisation-enforced, offline
and CPU-only, in the single-inference budget. The comparison with each family
makes the positioning concrete.

The generation-side copilots---PENTESTGPT
\parencite{deng2024pentestgpt} and the sparring-partner system
\parencite{happe2023pwned}---demonstrate that a model can drive security
tooling, but both route output to a human reviewer and validate correctness
informally. The IRE's checked output is structurally different: it is a
machine-consumed intent contract that has already passed deterministic
validation before any executor sees it. The measured 72\%/64\% catch rate does
not compete with those systems' task coverage; it answers a question they do
not pose, namely what fraction of flawed machine-bound decisions is caught
before dispatch.

\begin{sloppypar}
Fabrication-detection and agreement methods---SelfCheckGPT
\parencite{manakul2023selfcheckgpt}, semantic-uncertainty and
semantic-entropy approaches \parencite{kuhn2023semantic,farquhar2024semantic},
and self-consistency decoding \parencite{wang2023selfconsistency}---assume
repeated or interactive generation whose cost collides with the offline,
CPU-only, single-inference budget adopted here
(Section~\ref{sec:domain-foundations}). The IRE does not abandon their
rationale; it relocates it. Instead of sampling a verifier, it resolves the
cost tension with a bounded-sampling agreement gate and a deterministic
escalation path (the hallucination validator) that needs no large model
(Section~\ref{sec:related-research}). The measured greed of the semantic cache
(Section~\ref{sec:exp3}) is the accuracy-pressure relief that makes this
resolution practical on the constrained device.
\end{sloppypar}

The comparison with OWASP-prioritised defences
\parencite{owasp2025} and the empirically demonstrated insufficiency of
individual prevention techniques \parencite{liu2023promptinjection,
greshake2023indirect, kang2024dualuse} is the strongest one. Those surveyed
defences act as post-hoc application-layer checks; the IRE's contribution is
that the checks are composed into a stage that sits \emph{between} generation
and dispatch, where each predicate (schema, syntax, network, scope,
authorisation) is independently decidable on a closed intent taxonomy. The
ablation evidence that the complete stack outperforms any single layer is the
empirical counterpart of that architectural difference, and it is the reason
the comparison proceeds at the level of a composed system rather than of
individual techniques.

\section{Research Contribution}\label{sec:discussion-research}

The research contribution is the answered question, not the artefact. Three
stated findings constitute it.

\emph{First}, the deterministic validation stack is the dominant driver of the
measured safety, and its effect is reproducible across the two backends. The
ablation evidence in Section~\ref{sec:exp3} shows the stack removing
$30/24$ catch points against a $9/14$ contribution from the largest single
layer; the same ordering holds across backends. This generaliates the narrow
claim ``guardrails help'' into a measurable ordering result: which layers
matter, and in which combination.

\emph{Second}, the effectiveness boundary is predictable: catch concentrates
on flaw classes that expose decidable structure. The per-class catch spectrum
of Section~\ref{sec:exp4} runs from 100\% down to 40\%, and the ranking is
exactly the decidability ranking developed in the conceptual framework
(Section~\ref{sec:gap-synthesis}). This turns the limitation of single-shot
offline inference into a design law for trustworthy automation: a model in the
closed-intent setting should be asked for structure, never for verdicts.

\emph{Third}, usability and safety can improve together. The guarded pipeline
simultaneously raises accuracy, raises catch and halves false positives
(Section~\ref{sec:exp1}). The experiment record therefore weakens the frequent
implicit assumption that adding deterministic checks to model output must trade
either precision for recall or fluency for utility.

\section{Product Contribution}\label{sec:discussion-product}

The product contribution is the shippable IRE implementation developed in
Chapter~\ref{ch:implementation-testing} and exercised by the integration
analysis of Section~\ref{sec:integration-results}: a guarded intent-recognition
engine that converts free-form offensive-security instructions into validated,
schema-constrained, authorised machine contracts end-to-end. Its measured and
verifiable properties are an automated test suite of 506 tests (76\%
statement coverage, 434 unit and 72 integration) that encode the safety
invariant as executable expectations, and a live C1--C2 contract verification
that exercised the real code paths (Section~\ref{sec:integration-results}).

The product's differentiated value sits in the stack, not in the model: the
engine's own documentation and configuration define OWASP-aligned guardrail
hardening, and the safety gate rejects or warns on every flaw class rather
than rendering model output downstream. Within the deployment constraints--
offline, CPU-only, single-inference, retrieval-free--the product is the only
surveyed design in Table~\ref{tab:existing-systems} that provides the full
conjunction of schema check, scope/authorisation enforcement and offline
operation.

The practical caveats are the ones verified in the source audit
(Section~\ref{sec:requirements} and the integration record): several documented
dependencies are not implemented (flagged as \emph{doc-only}); the shipped
configuration uses a default development API key that bypasses transport
authentication when left untouched; and the ``\texttt{production}''
scope-mode value admitted by the settings schema is dead configuration that
falls through to the \texttt{WARN} enforcement path rather than
\texttt{BLOCK}. These are real product-integrity gaps and are stated as such;
none affects the validity of the evaluation, which ran with production
credentials and strict scope semantics, but each must be closed before the
engine is deployed in an untrusted environment.

\section{Limitations and Threats to Validity}\label{sec:discussion-limitations}

The limitations are stated precisely because the dissertation commits to
evidence over assertion. They are organised by validity category and, where
possible, bounded by the reproducibility record of
Section~\ref{sec:reproducibility}.

\emph{Single-run point estimates and the absence of confidence intervals}
are the principal internal-validity limitation. Each configuration in
Experiments~1--4 was executed once (Section~\ref{sec:procedure}), so no
measurement error or seed sweep bounds the deltas; the magnitude of individual
layer effects, especially the small negative-$\Delta$Catch entries, must be
read as one-sample observations. The protocol mitigates this by committing the
execution in advance and fixing the data set, but it does not remove it.

\emph{Synthetic, non-uniform data} bounds content validity. The 1,500-record
population is generated against a committed taxonomy, and its category sizes
are non-uniform (600 D+ / 900 D-, and per-class sizes that differ), so accuracy
and catch figures are population-specific and should not be read as
population-mean estimates for real production traffic.

\emph{Closed threat boundary} bounds external validity. The design explicitly
excludes retrieval and free-form rendering, which truncates the indirect- and
tool-injection channels, but the evaluation does not cover threats outside the
trust boundary: indirect prompt injection, physical device compromise, or
backdoored model weights were never presented to the engine. The claim is
therefore ``reliable within the evaluated threat model'', not ``reliable
unconditionally''.

\emph{Hardware-dependence.} All experiments ran on a single CPU-only device,
so conclusions about throughput, cache behaviour and feasibility are
single-platform observations. The hardware bound is recorded in the
reproducibility manifest (Section~\ref{sec:reproducibility}).

\emph{Unconfirmed mechanism.} The mechanism behind the
increase-on-removal pattern of the ablation's negative-$\Delta$Catch rows is
unconfirmed; it is attributed to removal-adversity interactions but no
mechanistic study isolates it. This is flagged as open rather than explained.

\emph{Weak protection of semantic classes.} The three semantic flaw classes
(contradictory, out-of-scope, ambiguous) are caught at 40--59\% because no
deterministic signal exists for them. The research question of how to improve
these classes within the offline budget is left open.

\emph{Engineering configuration caveats.} As stated in
Section~\ref{sec:discussion-product}, the default development API key and the
dead ``\texttt{production}'' scope mode are verified findings that materially
affect deployment posture and are recorded as product limitations, not
conflated with the evaluation.

\emph{Latency, energy and operational cost} are acknowledged as unmeasured
throughout. The evaluation prioritised correctness and catch; an operator
deploying the engine should measure end-to-end latency on their own hardware
before committing to interactive use.

\section{Implications}\label{sec:discussion-implications}

\subsection*{Implications for practice}

For practitioners deploying language models in security automation, the
strongest implication is architectural: put a deterministic, decidable
validation stage between a small offline model and machine dispatch, and the
safety and usability of the system can improve together rather than trade. The
catch-before-dispatch posture means the model is consulted for structure while
the system decides; this is directly actionable with offline inference and a
closed intent taxonomy, and it does not require a large verifier model. The
verified product gaps---default development credentials and the dead
\texttt{production} scope route---should be treated as design requirements for
any similar engine before deployment.

\subsection*{Implications for research}

The measured ordering result answers RQ3's conditional reading: safety
generalises across backends, accuracy tuning does not. That split suggests the
research frontier is not ``which model'' but ``which decidable predicate
structure'', and it strengthens the case for composing model behaviour with
enumerable checks over sampling-based agreement for single-inference budgets.
The weak-protection gap for the three semantic classes is the clearest open
problem: closing it within the offline budget (for example with stronger
structured prompting, targeted enumerators, or the bounded-sampling escalator)
is a concrete direction for future work. The single-run limitation should be
removed by repeating the committed protocol with seed sweeps and interval
estimates to convert the point estimates of this dissertation into bounded
ones.

\subsection*{Implications for the thesis}

For the dissertation itself, the implication is that the central question is
answered with the evidence available and the gaps are stated with the sources
that would close them. The evaluation remains valid within the closed threat
model and the committed protocol; the follow-up work---production provision
of \texttt{gemma4:e4b} for the C1--C2 happy-path re-run recorded in
Section~\ref{sec:integration-results}, latency measurement, and the semantic
class research---is itemised in the project's reproducibility record rather
than asserted.